\scanpage{87}
\mspage{36}

est dans $W$, de sorte que le diagramme
\diagram{\begin{tikzcd}[ampersand replacement=\&]
X'_0\arrow[r,"i_0","{\deleted{\unsure{W}}}"'] \& X' \& X'_1\arrow[l,"i_1"',"W"]
\end{tikzcd}}
définit une flèche \inserted{$\lambda(u)\colon$} $X'_0\to X'_1$\unsure{,} dans $\mathrm{Hot}_W$, savoir $\lambda(u)=i_1^{-1}\,i_0$. L'hyp. qui assure l'existence de $\lambda^\circ$ s'exprime par la condition symétrique $i_0\in W$, qui permet de définir la flèche $\lambda^\circ(u)\colon X'_1\mathbin{\unsure{\text{---}|\to}}X'_0$ par $\lambda^\circ(u)=i_0^{-1}i_1$.

Quand les deux conditions sont satisfaites simultanément, on voit donc que $\lambda(u),\lambda^\circ(u)$ \unsure{sont} inverses \unsure{l'un|l'une} de l'autre.

\uline{Remarques} Pour définir $\lambda$ ``universellement'', il n'était pas indispensable de considérer les $X_{/s}$ et les inclusions $X_s\mathbin{\text{---}}X_{/s}$, au lieu de cela il suffisait de regarder les chgts. de base $S'=\mathbb{\Delta}^1\xrightarrow{u}S$, et de faire l'hyp. que, pour tout tel chg. de base, l'inclusion
\[
X'_1\xrightarrow{i_1}X'=X_{(\mathbb{\Delta}^{1},u)}
\]
est dans $W$. Cette hypothèse est automatiquement stable par chg. de base arbitraire. (La $W$-propreté de $f\colon X\to S$ signifie que $X'_1\mathbin{\text{---}}X'$ est non seulement dans $W$, mais $\unsure{=}$ $W$-coasphérique.) Dualement, pour la définition ``universelle'' de $\lambda^\circ$, il \unsure{suffirait|suffisait} que pour tt chgt de base $\mathbb{\Delta}^1\to S$ comme dessus, l'inclusion $X'_0\xrightarrow{i_0}X'$ est dans $W$.

% archive sheet number: 87
